% steps/step07.tex
\section[گام هفتم: آموزش مدل]{\faChalkboardTeacher\ \ گام هفتم: آموزش مدل}
\begin{stepbox}

\field{انتخاب تابع زیان (\lr{Loss Function})}{
    در تصاویر \lr{CT}، ناحیه تومور معمولاً نسبت بسیار کوچکی از کل تصویر را تشکیل می‌دهد؛ بنابراین با عدم‌تعادل شدید کلاس‌ها مواجه هستیم و توابع زیان کلاسیک به‌تنهایی عملکرد مناسبی ندارند. رویکرد رایج و مؤثر، استفاده از \lr{Dice Loss} (برای بیشینه‌سازی همپوشانی ماسک پیش‌بینی‌شده و ماسک واقعی) در ترکیب با \lr{BCE} (\lr{Binary Cross-Entropy}) یا استفاده از \lr{Focal Loss} برای تمرکز بیشتر روی نمونه‌های سخت و کلاس اقلیت است. ضریب دایس به صورت زیر تعریف می‌شود:
    $$
        \mathrm{Dice}=\frac{2|X\cap Y|}{|X|+|Y|}
    $$
    که در آن $X$ ماسک پیش‌بینی‌شده و $Y$ ماسک \lr{Ground Truth} است.
}

\field{تنظیم ابرپارامترها (\lr{Hyperparameter Tuning})}{
    به دلیل حجم بالای داده‌های سه‌بعدی \lr{NIfTI} و سنگینی معماری \lr{UNet++} (یا \lr{Cascaded U-Net})، محدودیت حافظه \lr{GPU} تعیین‌کننده است؛ بنابراین \lr{Batch Size} معمولاً کوچک انتخاب می‌شود (حدود ۴ تا ۱۶). نرخ یادگیری (\lr{Learning Rate}) اغلب با مقدار اولیه $10^{-4}$ شروع می‌شود و استفاده از زمان‌بند نرخ یادگیری (\lr{LR Scheduler}) مانند \lr{ReduceLROnPlateau} برای کاهش خودکار نرخ در صورت توقف بهبود، توصیه می‌گردد. به‌عنوان بهینه‌ساز نیز الگوریتم‌های \lr{Adam} یا \lr{AdamW} انتخاب‌های استاندارد و کارآمدی برای این نوع شبکه‌ها هستند.
}

\newpage
\field{پایش فرآیند آموزش و جلوگیری از بیش‌برازش (\lr{Overfitting Prevention})}{
    به دلیل محدود بودن داده‌های برچسب‌خورده پزشکی، افزایش داده (\lr{Data Augmentation}) نقش کلیدی در بهبود تعمیم‌پذیری دارد. اعمال تبدیلاتی مانند چرخش (\lr{Rotation})، قرینه‌سازی (\lr{Flipping}) و به‌ویژه \lr{Elastic Deformation} (برای شبیه‌سازی تغییرات طبیعی بافت‌های نرم) بسیار مؤثر است. علاوه بر این، استفاده از توقف زودهنگام (\lr{Early Stopping}) ضروری است؛ به‌طوری که معیار \lr{Validation Dice Score} به‌صورت پیوسته پایش شود و اگر برای مثال طی ۱۰ ایپاک متوالی بهبودی مشاهده نشد، آموزش متوقف گردد تا از حفظیات روی داده‌های آموزش و افت عملکرد روی داده‌های جدید جلوگیری شود.
}

\end{stepbox}
